\chapter{因果效应的图形表示}

因果推断通常需要专家知识，以及一批关于连接处理、结局与其他变量的因果网络的无法检验的假设。此前各章关注的是在过分简化的场景中计算因果效应的条件与方法（例如，你抬头张望对其他行人行为的影响，一个理想化的心脏移植研究）。这样做的目的，是对真实情境所需的更精细方法背后的思想作一平易的引入。由于所考虑的场景如此简单，其实并不需要把因果网络明确写出来。然而，当我们开始把注意力转向更复杂的情境时，把关于相关变量知道什么、假设什么讲清楚，就变得至关重要。

本章引入一种图形工具，用来表示我们关于所关心的因果结构的定性专家知识与先验假设。图以直观的方式概括知识与假设，有助于厘清概念问题，也便于研究者之间交流。在因果推断问题中使用图，可以让你更容易听从一条明智的建议：先画出你的假设，再得出你的结论。

\section{因果图}

本章介绍图（graph），我们称之为\emph{因果图}（causal diagram），用以表示关键的因果概念。因果推断的现代图示理论产生于计算机科学与人工智能学科。

\mnote{关于这一主题的综合性著作见 Pearl (2009) 以及 Spirtes、Glymour 与 Scheines (2000)。}

本章及随后三章都围绕用因果图进行问题的概念化展开。

请看图 6.1 中的图。它包含三个节点，代表随机变量（$L$、$A$、$Y$），以及三条边（箭头）。我们采用时间从左向右流动的约定，因此 $L$ 在时间上先于 $A$ 与 $Y$，$A$ 在时间上先于 $Y$。与此前各章一样，$L$、$A$、$Y$ 分别代表疾病严重程度、心脏移植与死亡。

\begin{figure}[htbp]
\centering
\begin{tikzpicture}[causal]
  \node[vtx] (f1L) at (0,0) {$L$};
  \node[vtx] (f1A) at (1.7,0) {$A$};
  \node[vtx] (f1Y) at (3.4,0) {$Y$};
  \draw[ed] (f1L) -- (f1A);
  \draw[ed] (f1A) -- (f1Y);
\end{tikzpicture}
\caption{}
\label{fig:6.1}
\end{figure}

从某一变量 $V$ 指向另一变量 $W$ 的箭头存在，表示我们知道：至少对某一个体而言存在直接因果效应（即不经过图中任何其他变量中介的效应）。反之，箭头的缺失意味着我们知道：对总体中任何个体，$V$ 对 $W$ 都没有直接因果效应。例如在图 6.1 中，从 $L$ 指向 $A$ 的箭头意味着疾病严重程度影响接受心脏移植的概率。标准的因果图并不区分箭头代表的是有害效应还是保护效应。此外，若如图 6.1 那样，一个变量（这里是 $Y$）有两个原因，图也不编码这两个原因如何交互。

图 6.1 这类因果图称为\emph{有向无环图}（directed acyclic graph），通常缩写为 DAG。“有向”是因为边含有方向：从 $L$ 指向 $A$ 的箭头指向 $A$，故 $L$ 可能引起 $A$，反之则不然。“无环”是因为没有回路：一个变量不能直接或经由另一变量引起自身。

有向无环图还有因果推断之外的用途。这里我们关注的是\emph{因果}有向无环图。因果 DAG 的一个定义性性质是：以任一变量的直接原因为条件，图中该变量与任何它不是其原因的变量独立。这一假设称为\emph{因果马尔可夫假设}（causal Markov assumption），它蕴含：在因果 DAG 中，图中任意一对变量的共同原因也必须出现在图中。因果 DAG 的形式定义见技术注记 6.1。

\begin{techpoint}[因果有向无环图]
我们定义\emph{有向无环图}（DAG）$G$ 为一个图，其节点（顶点）是随机变量 $V=(V_1,\ldots,V_M)$，带有有向边（箭头），且没有有向回路。我们用 $\mathrm{PA}_m$ 表示 $V_m$ 的\emph{父节点}，即存在直接箭头指向 $V_m$ 的那些节点构成的集合。若存在一串由边连接、从 $V_j$ 出发沿箭头方向可到达 $V_m$ 的节点序列，则称 $V_m$ 是 $V_j$ 的\emph{后代}（而 $V_j$ 是 $V_m$ 的\emph{祖先}）。例如，考虑图 6.1 中的 DAG。该 DAG 中 $M=3$，可取 $V_1=L$、$V_2=A$、$V_3=Y$；$V_3=Y$ 的父节点 $\mathrm{PA}_3$ 为 $(L,A)$。我们采用如下排序约定：若 $m>j$，则 $V_m$ 不是 $V_j$ 的祖先。若对每个 $j$，$V_j$ 以其父节点为条件独立于其非后代，则称 $V$ 的分布关于 DAG $G$ 是\emph{马尔可夫的}（等价地，该分布按 DAG $G$ 分解）。后一陈述在数学上等价于：DAG $G$ 中变量 $V$ 的密度 $f(V)$ 满足马尔可夫分解
\[
f(v)=\prod_{j=1}^{M} f(v_j\mid \mathrm{pa}_j).
\]

\emph{因果 DAG} 是满足以下三条的 DAG：1）从节点 $V_j$ 到 $V_m$ 的箭头缺失（即 $V_j$ 不是 $V_m$ 的父节点）可解释为：相对于图中其他变量，$V_j$ 对 $V_m$ 没有直接因果效应；2）图中任意一对变量的所有共同原因，即便未被测量，本身也在图上；3）任何变量都是其后代的原因。除非我们作出一个把 DAG 所表示的因果结构与研究所得数据联系起来的假设，因果 DAG 没有实际用处。这一假设称为\emph{因果马尔可夫假设}，它断言：以直接原因为条件，变量 $V_j$ 与任何它不是其原因的变量独立。也就是说，以其父节点为条件，$V_j$ 独立于其非后代；因此，因果 DAG 关于 DAG $G$ 是马尔可夫的。
\end{techpoint}

例如，假设在我们的研究中，个体被随机分配到心脏移植 $A$，分配概率取决于其疾病严重程度 $L$。那么 $L$ 是 $A$ 与 $Y$ 的共同原因，需要纳入图中，如图 6.1 的因果图所示。现在假设在我们的研究中，所有个体被随机分配到心脏移植的概率相同，与疾病严重程度无关。那么 $L$ 不是 $A$ 与 $Y$ 的共同原因，不必纳入因果图。图 6.1 表示\emph{条件}随机化实验，而图 6.2 表示\emph{边缘}随机化实验。

\begin{figure}[htbp]
\centering
\begin{tikzpicture}[causal]
  \node[vtx] (f2A) at (0,0) {$A$};
  \node[vtx] (f2Y) at (1.9,0) {$Y$};
  \draw[ed] (f2A) -- (f2Y);
\end{tikzpicture}
\caption{}
\label{fig:6.2}
\end{figure}

图 6.1 也可以表示一项观察性研究。具体而言，图 6.1 所表示的观察性研究中，我们愿意假设心脏移植 $A$ 的分配只有疾病严重程度 $L$ 这一个原因（父节点），且 $Y$ 没有其他原因。否则，那些 $Y$ 的原因即便未被测量，也需要纳入图中，因为它们会是 $A$ 与 $Y$ 的共同原因。下一章我们将说明：愿意把图 6.1 当作某项观察性研究的因果图，正是“给定 $L$ 的条件可交换性”这一假设——即对所有 $a$ 有 $Y^{a}\indep A\given L$——的图形翻译。

许多人觉得，用图形方法做因果推断比用反事实方法更易上手、更直观。然而这两种方法联系紧密。具体而言，每张图都对应一个底层的反事实模型（见技术注记 6.2 与 6.3）。正是这个模型为我们下面描述的启发式、直观的图形方法提供了数学依据。不过，常规因果图并不把底层的反事实变量画在图上。因此，图与反事实之间的联系历来是隐而不显的。最近发展起来的一类因果有向无环图——\emph{单一世界干预图}（Single World Intervention Graph，SWIG）——把反事实变量明确画在图上，从而无缝地统一了因果推断的反事实方法与图形方法。我们把 SWIG 的引入推迟到第 7 章，因为本章的内容是必要的预备。

\mnote{Richardson 与 Robins (2013) 提出了单一世界干预图（SWIG）。}

\begin{techpoint}[与因果 DAG 相联系的反事实模型]
在本书中，因果 DAG $G$ 表示一个底层的反事实模型。为给出 DAG $G$ 所表示的反事实模型的形式定义，我们使用如下记号。对任意随机变量 $W$，记 $\mathcal{W}$ 为 $W$ 的支撑集（即 $W$ 的一切可能取值 $w$ 的集合）。对任意一组有序变量 $W_1,\ldots,W_m$，定义 $\bar{w}_m=(w_1,\ldots,w_m)$。设 $R$ 为 $V$ 的任一子集，$r$ 为 $R$ 的一个取值。则 $V_m^{r}$ 表示把 $R$ 置为 $r$ 时 $V_m$ 的反事实取值。

由顶点集 $V=(V_1,V_2,\ldots,V_M)$ 的 DAG $G$ 所表示的\emph{非参数结构方程模型}（NPSEM）（顶点按“$i<j$ 时 $V_i$ 不是 $V_j$ 的后代”排序）假定存在不可观测随机变量（误差）$\epsilon_m$ 与确定性未知函数 $f_m(\mathrm{pa}_m,\epsilon_m)$，使得 $V_1=f_1(\epsilon_1)$，且一步前向反事实 $V_m^{v_{m-1}}=f_m(\mathrm{pa}_m,\epsilon_m)$。也就是说，相对于 $G$ 上的其他变量，只有 $V_m$ 的父节点对 $V_m$ 有直接效应。NPSEM 蕴含图上任何变量都可被干预，因为假定把 $V_j$ 置为特定值 $v_j$ 的反事实存在。事实变量 $V_m$ 与对任意 $R\subset V$ 的反事实 $V_m^{r}$ 都可由 $V_1$ 与 $V_j^{v_{j-1}}$（$M\geq j>1$）递归得到。例如，$V_3^{v_1}=V_3^{v_1,v_2}$，即当 $V_1$ 被置为 $v_1$ 时 $V_3$ 的反事实取值 $V_3^{v_1}$，就是一步前向反事实 $V_3^{v_1,v_2}$，其中 $v_2$ 等于 $V_2$ 的反事实取值 $V_2^{v_1}$。类似地，$V_3=V_3^{v_1,V_2^{v_1}}$，且 $V_3^{v_1,v_4}=V_3^{v_1}$，因为 $V_4$ 不是 $V_3$ 的直接原因。从 $V_j$ 到 $V_k$ 的箭头缺失蕴含：对任何个体，$V_j$ 都不是 $V_k$ 的直接原因。

Robins (1986) 引入这一 NPSEM，称之为“与数据同样详细的”\emph{最精细因果解释结构树图}（finest causally interpreted structural tree graph，FCISTG），并把 $V_m$ 的父节点 $\mathrm{PA}_m$ 称为在控制排序中较早变量后 $V_m$ 的因果风险因子。Pearl (2009) 说明如何用 DAG 表示这一模型。Robins (1986) 还提出了往往更切实际的因果解释结构树图，其中只有一部分变量可被干预。为阐述方便，我们一般假定每个变量都可被干预，尽管这里考虑的统计方法实际上并不需要这一假设。
\end{techpoint}

\begin{techpoint}[与反事实模型相联系的独立性]
FCISTG 模型并不蕴含技术注记 6.1 的因果马尔可夫假设成立；还需要额外的统计独立假设。例如，Pearl (2000) 通常假定一个 NPSEM，其中所有误差项 $\epsilon_m$ 相互独立。我们把 Pearl 的误差独立模型称为 NPSEM-IE。相比之下，Robins (1986) 只假定：给定任意 $v_M$，一步前向反事实 $V_m^{v_{m-1}}=f_m(\mathrm{pa}_m,\epsilon_m)$（$m=1,\ldots,M$）联合独立，其中 $v_{m-1}$ 是 $v_M$ 的子向量；他称之为“与数据同样详细的”\emph{最精细全随机化因果解释结构树图}（finest fully randomized causally interpreted structured tree graph，FFRCISTG）模型。

更精确地说，Robins (1986) 假定：对每个 $m$，以事实过去 $V^{m-1}=v^{m-1}$ 为条件，从 $m+1$ 起的任何未来演化（与 $v^{m-1}$ 相容的一步前向反事实）都独立于事实变量 $V_m$。Robins 与 Richardson (2010) 证明，对正值分布，这一假设等价于上一段的假设。在无正值性时，我们按上一段定义模型。

Robins (1986) 证明，他的独立假设蕴含因果马尔可夫假设成立。NPSEM-IE 是 FFRCISTG，但反之不成立，因为 NPSEM-IE 作出的独立假设远多于 FFRCISTG (Robins and Richardson 2010)。

除另有说明，DAG 表示一个 NPSEM，但可能需指明是哪种类型。例如，图 6.2 的 DAG 既可以对应蕴含完全可交换性 $(Y^{a=0},Y^{a=1})\indep A$ 的 NPSEM-IE，也可以对应只蕴含 $a=0$ 与 $a=1$ 时边缘可交换性 $Y^{a}\indep A$ 的 FFRCISTG。凡未提及底层模型时，我们假定因果 DAG 表示与数据同样详细的 FFRCISTG。
\end{techpoint}

因果图是编码我们关于问题定性因果结构的主题知识与假设的简便方式。但如下文各节所述，因果图也编码因果网络中变量之间潜在关联的信息。正是这种对关联与因果的同时表示，使因果图成为如此有吸引力的工具。下文对从因果图推断关联的图形规则作一非正式介绍。我们强调概念上的洞见，而非形式上的严格。

\section{因果图与边缘独立}

考虑下面两个例子。第一，假设你知道阿司匹林使用 $A$ 对心脏病风险 $Y$ 有预防性因果效应，即 $\Pr[Y^{a=1}=1]\neq\Pr[Y^{a=0}=1]$。在一个阿司匹林 $A$ 被随机、无条件分配的实验中，图 6.2 的因果图就是这一知识的图形翻译。第二，假设你知道随身携带打火机 $A$ 对任何人患肺癌的风险 $Y$ 都没有因果效应（无论致病还是防病），即 $\Pr[Y^{a=1}=1]=\Pr[Y^{a=0}=1]$，而吸烟 $L$ 对随身携带打火机 $A$ 与肺癌 $Y$ 都有因果效应。图 6.3 的因果图就是这一知识的图形翻译。$A$ 与 $Y$ 之间没有箭头，表示携带打火机对肺癌没有因果效应；$L$ 被画成 $A$ 与 $Y$ 的共同原因。

\begin{figure}[htbp]
\centering
\begin{tikzpicture}[causal]
  \node[vtx] (f3L) at (0,0) {$L$};
  \node[vtx] (f3A) at (2.3,0.75) {$A$};
  \node[vtx] (f3Y) at (2.3,-0.75) {$Y$};
  \draw[ed] (f3L) -- (f3A);
  \draw[ed] (f3L) -- (f3Y);
\end{tikzpicture}
\caption{}
\label{fig:6.3}
\end{figure}

画图 6.2 与 6.3 时，我们只用了你关于图中变量间因果关系的知识；但有趣的是，这些因果图也编码了图中变量之间预期关联（更确切地说，预期缺乏关联）的信息。下面我们启发式地论证：一般而言，在图 6.2 与 6.3 中 $A$ 与 $Y$ 都会有关联；并叙述因果图理论中若干关键的相关结果。

先看图 6.2 所表示的随机化实验。直觉上，人们会预期仅由一条因果箭头连接的两个变量 $A$ 与 $Y$ 有关联。因果图理论恰恰表明了这一点：当已知 $A$ 对 $Y$ 有因果效应时（如图 6.2），一般也应预期 $A$ 与 $Y$ 有关联。这当然与下述事实一致：在具有无条件可交换性的理想随机化实验中，因果 $\Pr[Y^{a=1}=1]\neq\Pr[Y^{a=0}=1]$ 蕴含关联 $\Pr[Y=1\given A=1]\neq\Pr[Y=1\given A=0]$，反之亦然。捕捉因果图中因果—关联对应关系的一个启发式想法，是把两变量之间的路径想象成输送关联的管道或导线。

\mnote{DAG 中两个变量 $R$ 与 $S$ 之间的\emph{路径}，是沿一串边连接 $R$ 与 $S$、且途中不重复访问任何变量的路线。若一条路径完全由箭头方向相同的边组成，则称它为\emph{因果路径}，否则称为\emph{非因果路径}。}

关联与因果不同，是两个变量之间的对称关系；因此，关联一旦存在，就在两变量之间流动，与因果箭头的方向无关。在图 6.2 中，你也可以等价地说关联从 $A$ 流向 $Y$，或从 $Y$ 流向 $A$。

现在考虑图 6.3 所表示的观察性研究。我们知道携带打火机 $A$ 对肺癌 $Y$ 没有因果效应。现在的问题是：携带打火机 $A$ 是否与肺癌 $Y$ 有关联？我们知道 $\Pr[Y^{a=1}=1]=\Pr[Y^{a=0}=1]$，但是否也有\mbox{$\Pr[Y=1\given A=1]=$}\mbox{$\Pr[Y=1\given A=0]$}？要回答这个问题，设想一位天真的研究者决定研究携带打火机 $A$ 对肺癌风险 $Y$ 的影响（我们知道没有影响，但研究者并不知道）。他询问一大批人是否随身携带打火机，然后记录他们在随后 5 年内是否被诊断出肺癌。赫拉是研究参与者之一。我们得知赫拉随身携带打火机。但如果赫拉携带打火机（$A=1$），她更可能是一名吸烟者（$L=1$），因此她患肺癌（$Y=1$）的风险高于平均水平。于是我们直觉地得出结论：预期 $A$ 与 $Y$ 有关联，因为携带打火机者（$A=1$）的癌症风险不同于不携带者（$A=0$），即 $\Pr[Y=1\given A=1]\neq\Pr[Y=1\given A=0]$。换言之，即便 $A$ 对 $Y$ 没有因果效应，掌握关于处理 $A$ 的信息仍能提高我们预测结局 $Y$ 的能力。如果研究者仅因为 $A$ 与 $Y$ 有关联就断定 $A$ 对 $Y$ 有因果效应，他就错了。因果图理论再次证实了我们的直觉。用图形语言说，$A$ 与 $Y$ 有关联，是因为存在一条经由共同原因 $L$ 从 $A$ 到 $Y$（等价地，从 $Y$ 到 $A$）的关联流。

再看第三个例子。假设你知道某种基因单倍型 $A$ 对任何人成为吸烟者 $Y$ 的风险都没有因果效应，即 $\Pr[Y^{a=1}=1]=\Pr[Y^{a=0}=1]$，而单倍型 $A$ 与吸烟 $Y$ 对心脏病风险 $L$ 都有因果效应。图 6.4 的因果图就是这一知识的图形翻译。$A$ 与 $Y$ 之间没有箭头，表示该单倍型对吸烟没有因果效应；$L$ 被画成 $A$ 与 $Y$ 的共同效应。共同效应 $L$ 称为路径 $A\to L\leftarrow Y$ 上的\emph{对撞节点}（collider），因为两个箭头在该节点上相撞。

\begin{figure}[htbp]
\centering
\begin{tikzpicture}[causal]
  \node[vtx] (f4A) at (0,0.85) {$A$};
  \node[vtx] (f4Y) at (0,-0.85) {$Y$};
  \node[vtx] (f4L) at (2.3,0) {$L$};
  \draw[ed] (f4A) -- (f4L);
  \draw[ed] (f4Y) -- (f4L);
\end{tikzpicture}
\caption{}
\label{fig:6.4}
\end{figure}

问题同样是 $A$ 与 $Y$ 是否有关联。要回答它，设想另一位研究者决定研究单倍型 $A$ 对成为吸烟者 $Y$ 的风险的影响（我们知道没有影响，但研究者并不知道）。她对一大批儿童做基因测定，然后记录他们最终是否成为吸烟者。阿波罗是研究参与者之一。我们得知阿波罗没有该单倍型（$A=0$）。他成为吸烟者（$Y=1$）的可能性比一般人更高还是更低？了解单倍型 $A$ 并不能提高我们预测结局 $Y$ 的能力，因为有（$A=1$）与没有（$A=0$）该单倍型者的风险相同，即 $\Pr[Y=1\given A=1]=\Pr[Y=1\given A=0]$。换言之，我们会直觉地得出结论：$A$ 与 $Y$ 没有关联，即 $A$ 与 $Y$ 独立，或 $A\indep Y$。$A$ 与 $Y$ 都引起心脏病 $L$ 这一知识，在考虑 $A$ 与 $Y$ 之间的关联时无关紧要。因果图理论再次证实我们的直觉，因为它指出：对撞节点与其他变量不同，会阻断其所在路径上的关联流。因此 $A$ 与 $Y$ 独立，因为二者之间唯一的路径 $A\to L\leftarrow Y$ 被对撞节点 $L$ 阻断。

总之，两个变量若一个引起另一个，或共享共同原因，则（边缘地）有关联；否则它们（边缘地）独立。下一节探讨两个变量 $A$ 与 $Y$ 在给定第三个变量 $L$ 的条件下可能独立的条件。

\section{因果图与条件独立}

现在我们回到图 6.2、6.3、6.4 所描绘的情境，讨论因果图中的条件独立概念。

按照图 6.2，我们预期阿司匹林 $A$ 与心脏病 $Y$ 有关联，因为阿司匹林对心脏病有因果效应。现在假设我们获得一条额外信息：阿司匹林 $A$ 之所以影响心脏病风险 $Y$，是因为它降低血小板聚集 $B$。这一新知识被翻译为图 6.5 的因果图，其中血小板聚集 $B$（1：高，0：低）是 $A$ 对 $Y$ 之效应的中介。

\begin{figure}[htbp]
\centering
\begin{tikzpicture}[causal]
  \node[vtx] (f5A) at (0,0) {$A$};
  \node[vtx] (f5B) at (2.1,0) {$B$};
  \node[vtx] (f5Y) at (4.2,0) {$Y$};
  \draw[ed] (f5A) -- (f5B);
  \draw[ed] (f5B) -- (f5Y);
  \node[draw,inner sep=2.6pt,line width=0.5pt,fit=(f5B)] {};
\end{tikzpicture}
\caption{}
\label{fig:6.5}
\end{figure}

一旦在因果图中引入第三个变量，就可以提出一个新问题：在 $B$ 的各水平内（以 $B$ 为条件），$A$ 与 $Y$ 之间是否有关联？或者等价地：当我们已掌握 $B$ 的信息时，关于 $A$ 的信息是否还能提高预测 $Y$ 的能力？要回答这个问题，假设在一大批个体中收集了 $A$、$B$、$Y$ 的数据，并把分析限制在血小板聚集低（$B=0$）的那部分个体。图 6.5 中节点 $B$ 周围的方框表示这一限制。（若分析限制在 $B=1$ 的那部分个体，我们同样会在 $B$ 周围画一个方框。）

\mnote{由于在完整因果图（所有可能的箭头都存在的图）中不预期任何条件独立，人们常说关联的信息藏在缺失的箭头里。}

血小板聚集低（$B=0$）的个体患心脏病的风险低于平均水平。现在取其中一位个体。无论该个体是受治疗（$A=1$）还是未受治疗（$A=0$），我们早已知道他因血小板聚集低而风险低于平均水平。事实上，由于阿司匹林使用只通过血小板聚集影响心脏病风险，了解一个个体的治疗状态并不能为预测其心脏病风险提供任何额外信息。因此，在 $B=0$ 的那部分个体中，处理 $A$ 与结局 $Y$ 没有关联。（对 $B=1$ 组的个体也可作同样的非正式论证。）尽管 $A$ 与 $Y$ 边缘有关联，但给定 $B$ 时 $A$ 与 $Y$ 条件独立（无关联），因为在 $B$ 的各水平内，受治疗者与未受治疗者的心脏病风险相同：对所有 $b$ 有 \mbox{$\Pr[Y=1\given A=1,B=b]=\Pr[Y=1\given A=0,B=b]$}。即 $A\indep Y\given B$。用图形语言说，围绕变量 $B$ 的方框阻断了通过路径 $A\to B\to Y$ 的关联流。

现在回到图 6.3。上一节我们得出结论：携带打火机 $A$ 与肺癌风险 $Y$ 有关联，因为路径 $A\leftarrow L\to Y$ 对从 $A$ 到 $Y$ 的关联流是开放的。现在要问：给定 $L$ 时 $A$ 是否与 $Y$ 有关联？这个新问题由图 6.6 中 $L$ 周围的方框表示。假设研究者把研究限制在非吸烟者（$L=0$）。在这种情况下，得知某个个体携带打火机（$A=1$）无助于预测其肺癌风险（$Y=1$），因为“预测更好”的整个论证依赖于携带打火机者更可能是吸烟者这一事实。当研究限制在非吸烟者，或更一般地限制在具有某一特定吸烟强度的人群时，这一论证便无关紧要。尽管 $A$ 与 $Y$ 边缘有关联，但给定 $L$ 时 $A$ 与 $Y$ 条件独立，因为在 $L$ 的各水平内，受治疗者与未受治疗者的肺癌风险相同：对所有 $l$ 有 \mbox{$\Pr[Y=1\given A=1,L=l]=\Pr[Y=1\given A=0,L=l]$}。即 $A\indep Y\given L$。用图形语言说，$A$ 与 $Y$ 之间的关联流被打断，因为路径 $A\leftarrow L\to Y$ 被 $L$ 周围的方框阻断。

\begin{figure}[htbp]
\centering
\begin{tikzpicture}[causal]
  \node[vtx] (f6L) at (0,0) {$L$};
  \node[vtx] (f6A) at (2.4,0.85) {$A$};
  \node[vtx] (f6Y) at (2.4,-0.85) {$Y$};
  \draw[ed] (f6L) -- (f6A);
  \draw[ed] (f6L) -- (f6Y);
  \node[draw,inner sep=2.6pt,line width=0.5pt,fit=(f6L)] {};
\end{tikzpicture}
\caption{}
\label{fig:6.6}
\end{figure}

\mnote{通过共同原因阻断处理与结局之间的关联流，就是以分层作为达到可交换性之方法的图形依据。}

最后，再考虑图 6.4。上一节我们得出结论：具有单倍型 $A$ 与成为吸烟者 $Y$ 独立，因为 $A$ 与 $Y$ 之间的路径 $A\to L\leftarrow Y$ 被对撞节点 $L$ 阻断。现在我们启发式地论证：一般而言，在共同效应 $L$ 的各水平内，$A$ 与 $Y$ 会条件关联。假设研究者想估计单倍型 $A$ 对吸烟状态 $Y$ 的效应，把研究总体限制为患心脏病者（$L=1$）。图 6.7 中 $L$ 周围的方框表示他们以 $L$ 的某一特定取值为条件。知道一个患心脏病的人没有单倍型 $A$，会提供一些关于其吸烟状态的信息，因为在没有 $A$ 的情况下，更可能存在 $L$ 的另一个原因，例如 $Y$。也就是说，在患心脏病者中，没有单倍型 $A$ 的人里吸烟者比例更高。因此，给定 $L=1$ 时，$A$ 与 $Y$ 呈反向关联。如果研究者仅因为 $A$ 与 $Y$ 在 $L$ 各水平内有关联就断定 $A$ 对 $Y$ 有因果效应，他就错了。极端情形下，若 $A$ 与 $Y$ 是 $L$ 仅有的两个原因，那么在患心脏病者中，缺少其中一个便能完美预测另一个的存在。因果图理论表明，以 $L$ 这样的对撞节点为条件确实打开了路径 $A\to L\leftarrow Y$，而该路径在未以对撞节点为条件时是阻断的。直觉上，两个变量（原因）是否有关联，不可能受一个未来事件（它们的效应）影响；但一旦我们按共同效应分层，某一效应的两个原因一般就会变得有关联。

\begin{figure}[htbp]
\centering
\begin{tikzpicture}[causal]
  \node[vtx] (f7A) at (0,0.85) {$A$};
  \node[vtx] (f7Y) at (0,-0.85) {$Y$};
  \node[vtx] (f7L) at (2.4,0) {$L$};
  \draw[ed] (f7A) -- (f7L);
  \draw[ed] (f7Y) -- (f7L);
  \node[draw,inner sep=2.6pt,line width=0.5pt,fit=(f7L)] {};
\end{tikzpicture}
\caption{}
\label{fig:6.7}
\end{figure}

再举一例：图 6.8 的因果图在图 6.7 的基础上增加了一种利尿药物 $C$，其使用是心脏病诊断的结果。在 $C$ 的各水平内，$A$ 与 $Y$ 也有关联，因为 $C$ 是 $A$ 与 $Y$ 的共同效应。因果图理论表明，以受对撞节点 $L$ 影响的变量 $C$ 为条件，同样会打开路径 $A\to L\leftarrow Y$。若不针对对撞节点 $L$ 或其后果 $C$ 作条件，这条路径是阻断的。

\begin{figure}[htbp]
\centering
\begin{tikzpicture}[causal]
  \node[vtx] (f8A) at (0,0.9) {$A$};
  \node[vtx] (f8Y) at (0,-0.9) {$Y$};
  \node[vtx] (f8L) at (2.3,0) {$L$};
  \node[vtx] (f8C) at (4.4,0) {$C$};
  \draw[ed] (f8A) -- (f8L);
  \draw[ed] (f8Y) -- (f8L);
  \draw[ed] (f8L) -- (f8C);
  \node[draw,inner sep=2.6pt,line width=0.5pt,fit=(f8C)] {};
\end{tikzpicture}
\caption{}
\label{fig:6.8}
\end{figure}

\mnote{以共同效应为条件所致关联的更多讨论，见第 8 章。}

本节与上一节考察了两个变量可能有关联的三种结构性原因：一个引起另一个；它们共享共同原因；它们共享共同效应，且分析被限制在该共同效应（或其后代）的某一水平。在此过程中，我们引入了若干可应用于任何因果图、用以判定两个变量是否（条件）独立的图形规则。我们用来支持这些图形规则的论证是启发式的，依赖于我们的因果直觉。然而这些论证已被形式化并得到数学证明。系统的图形规则总结见精要 6.1，忠实性概念的介绍见精要 6.2。

\mnote{这些图形规则背后的数学理论称为“d 分离”（Pearl 1995）。}

两个变量之间还有一种我们尚未讨论的可能关联来源：机遇或随机变异。与两变量关联的结构性原因——一个对另一个的因果效应、共享共同原因、以共同效应为条件——不同，随机变异导致的偶然关联会随着研究总体规模增大而变小。

为使讨论集中于结构性关联而非偶然关联，我们继续假定（直至第 10 章）已记录了一个非常大（也许是假设的）目标总体中每个个体的数据。

\begin{finepoint}[d 分离]
我们按如下图形规则把一条路径定义为\emph{阻断的}或\emph{开放的}。

\begin{enumerate}
\item 若没有对任何变量取条件，则一条路径是阻断的，当且仅当路径上有两个箭头在路径的某个变量上相撞。在图 6.1 中，路径 $L\to A\to Y$ 是开放的，而路径 $A\to Y\leftarrow L$ 是阻断的，因为路径上的两个箭头在 $Y$ 处相撞。我们称 $Y$ 是路径 $A\to Y\leftarrow L$ 上的对撞节点。
\item 任何包含已被取条件的非对撞节点的路径都是阻断的。在图 6.5 中，以 $B$ 为条件后，$A$ 与 $Y$ 之间的路径是阻断的。我们用变量周围的方框表示对它取条件。
\item 已被取条件的对撞节点不阻断路径。在图 6.7 中，以 $L$ 为条件后，$A$ 与 $Y$ 之间的路径是开放的。
\item 其后代已被取条件的对撞节点不阻断路径。在图 6.8 中，以对撞节点 $L$ 的后代 $C$ 为条件后，$A$ 与 $Y$ 之间的路径是开放的。
\end{enumerate}

规则 1--4 可概括如下。一条路径是阻断的，当且仅当它包含一个已被取条件的非对撞节点，或包含一个未被取条件、且没有已被取条件的后代的对撞节点。若两个变量之间的所有路径都被阻断，则称二者 \emph{d 分离}（否则称 \emph{d 连通}）。若第一集合中的每个变量都与第二集合中的每个变量 d 分离，则称两个变量集合 d 分离。因此，在图 6.1 中 $A$ 与 $L$ 不是 d 分离的，因为二者之间有一条开放路径（$L\to A$），尽管另一条路径（$A\to Y\leftarrow L$）被对撞节点 $Y$ 阻断。然而在图 6.4 中，$A$ 与 $Y$ 是 d 分离的，因为二者之间唯一的路径被对撞节点 $L$ 阻断。

统计独立与纯图形的 d 分离概念之间的关系依赖于因果马尔可夫假设（技术注记 6.1）：在因果 DAG 中，任何变量以其父节点为条件独立于其非后代。Pearl (1988) 证明了如下基本定理：因果马尔可夫假设蕴含，对任意三个互不相交的变量集合 $A$、$B$、$C$，若 $A$ 在给定 $C$ 时与 $B$ d 分离，则 $A$ 在给定 $C$ 时与 $B$ 统计独立。其逆命题也成立的假设，即若 $A$ 在给定 $C$ 时与 $B$ 统计独立则 $A$ 在给定 $C$ 时与 $B$ d 分离，是另一条独立假设——精要 6.2 所述的忠实性假设。在忠实性下，图 6.5 中 $A$ 在给定 $B$ 时与 $Y$ 条件独立，图 6.7 中 $A$ 在给定 $L$ 时与 $Y$ 不条件独立，图 6.8 中 $A$ 在给定 $C$ 时与 $Y$ 不条件独立。从因果图推断关联陈述的 d 分离规则（“d-”代表方向性）由 Pearl (1995) 形式化。一套等价的图形规则称为“道德化”（moralization），由 Lauritzen 等 (1990) 提出。
\end{finepoint}

\begin{finepoint}[忠实性]
在因果 DAG 中，从 $A$ 到 $Y$ 的箭头缺失表示“$A$ 对任何个体的 $Y$ 都没有因果效应”这一严格零假设成立；而箭头 $A\to Y$（如图 6.2）表示 $A$ 对总体中至少一个个体的结局 $Y$ 有因果效应。因此，我们一般会预期：在图 6.2 下，$A$ 对 $Y$ 的平均因果效应 $\Pr[Y^{a=1}=1]-\Pr[Y^{a=0}=1]$ 与 $A$ 和 $Y$ 之间的关联 $\Pr[Y=1\given A=1]-\Pr[Y=1\given A=0]$ 都不为零。然而这未必成立：由图 6.2 表示的情境，也可能是既无平均因果效应、又无关联的情境。例如，回想表 4.1 的数据。心脏移植 $A$ 增加女性（总体的一半）的死亡风险 $Y$，降低男性（另一半）的死亡风险。由于 $A$ 的有益效应与有害效应恰好相互抵消，平均因果效应为零，$\Pr[Y^{a=1}=1]=\Pr[Y^{a=0}=1]$。然而图 6.2 仍是正确的因果图，因为处理 $A$ 影响总体中某些个体——事实上是所有个体——的结局 $Y$。

形式上，忠实性是如下假设：对因果 DAG 上三个互不相交的集合 $A$、$B$、$C$（$C$ 可为空集），若给定 $C$ 时 $A$ 与 $B$ 独立，则给定 $C$ 时 $A$ 与 $B$ d 分离。当因果图使我们预期存在某种关联、而数据中实际并不存在时（如本例），我们说数据的联合分布对该因果 DAG \emph{不忠实}。在本例中，不忠实是效应修饰（按性别）造成的，两半总体中的效应方向相反而大小恰好相等。这种效应的完美抵消很罕见，因此本书通篇假定忠实性。由于不忠实的分布罕见，在实践中几乎总可以把缺乏 d 分离（见精要 6.1）等同于非零关联。

然而，有些情形中忠实性是被设计所破坏的。例如，考虑 4.5 节的前瞻性研究。$A$ 对 $Y$ 的平均因果效应是在按 $L$ 匹配后计算的。在匹配后的总体中，$L$ 与 $A$ 没有关联，因为 $L$ 的分布在受治疗者与未受治疗者中相同。也就是说，个体之所以被选入匹配总体，是因为他们具有 $L$ 与 $A$ 的某种特定取值组合。图 6.9 的因果图表示匹配研究的情境，其中选择 $S$（1：是，0：否）由 $A$ 与 $L$ 共同决定。$S$ 周围的方框表示分析限于被选入匹配队列者（$S=1$）。按 d 分离规则，以 $S$ 为条件时，$A$ 与 $L$ 之间有两条开放路径：$L\to A$ 与 $L\to S\leftarrow A$。因此人们会预期给定 $S$ 时 $L$ 与 $A$ 有关联。然而匹配确保 $L$ 与 $A$ 没有关联（见第 4 章）。为什么会有这种差异？匹配通过路径 $L\to S\leftarrow A$ 制造了一种关联，其大小与经由路径 $L\to A$ 的关联相等而方向相反。净结果是两种关联完美抵消。匹配导致不忠实。

\begin{center}
\begin{tikzpicture}[causal]
  \node[vtx] (f9L) at (0,0.95) {$L$};
  \node[vtx] (f9A) at (2.2,0.95) {$A$};
  \node[vtx] (f9Y) at (4.4,0.95) {$Y$};
  \node[vtx] (f9S) at (3.4,-0.75) {$S$};
  \draw[ed] (f9L) -- (f9A);
  \draw[ed] (f9A) -- (f9Y);
  \draw[ed] (f9L) to[bend left=32] (f9Y);
  \draw[ed] (f9L) -- (f9S);
  \draw[ed] (f9A) -- (f9S);
  \node[draw,inner sep=2.6pt,line width=0.5pt,fit=(f9S)] {};
\end{tikzpicture}
\captionof{figure}{}
\label{fig:6.9}
\end{center}

最后，当图上变量之间存在确定性关系时，忠实性也可能被破坏。具体而言，当两个变量由含有确定性箭头的路径相连时，若二者之间的所有路径都被阻断，则二者独立；但即便某些路径开放，二者也可能独立。本书除另有说明外假定忠实性。当数据分析的目标是发现因果结构时，也假定忠实性（见精要 6.3）。
\end{finepoint}

\section{因果图中的正值性与一致性}

由于因果图编码我们关于因果结构的定性专家知识，它们可用作直观辅助，帮助把因果问题概念化并指导数据分析。事实上，第 2 章所述用于量化处理效应的公式——标准化与逆概率加权——也可以用因果图理论推导，属于有时称为 do 演算（do-calculus）的一部分。因此，我们在第 1--5 章选择反事实理论，其实并未偏袒某一种特定方法，而只是偏袒某一种特定记号。

\mnote{Pearl (2009) 综述了由图论推导出的定量因果推断方法。}

无论使用哪种记号（反事实还是图），可交换性、正值性与一致性都是通过标准化或逆概率加权进行因果推断所需的条件。若其中任一条件不成立，数据分析得到的数字就可能无法恰当地解释为因果效应的度量。下一节（以及第 7、8 章）讨论可交换性条件如何翻译为图的语言。这里我们关注正值性与一致性。

遗憾的是，因果图无法编码对正值性的违背，除非在特殊情形下，例如正值性被违背是因为处理 $A$ 恰好是某个处理前变量 $L$ 的确定性函数。在这种情况下，我们把 $L\to A$ 箭头加粗以表示确定性。一致性的第一个组成部分——定义良好的反事实——意味着从处理 $A$ 到结局 $Y$ 的箭头对应一个可能假想但相对明确的干预（或者，当 $A$ 是一种状态时，意味着存在只通过 $A$ 影响 $Y$ 的定义良好的干预；见精要 6.4）。在本书讨论的因果图中，除另有说明，正值性是隐含的；一致性则嵌入记号之中，因为我们只考虑具有相对定义良好的干预的处理节点。注意，正值性关心的是指向处理节点的箭头，而定义良好的干预只关心从处理节点指出的箭头。

于是，处理节点被隐含地赋予了与其他所有节点不同的地位。一些作者通过在因果图中加入决策节点来明确这一差别。

\mnote{影响图（influence diagram）是加入决策节点、以表示所关心干预的因果图 (Dawid 2000, 2002)。}

尽管这种决策论取向在很大程度上导致与本书所述相同的方法，我们在本章给出的因果图中不加入决策节点。由于我们总是明确说明对变量 $A$ 的可能干预，额外的节点（用来表示可能干预）会有些冗余。不过，在后续各章使用 SWIG——其节点表示反事实变量的因果图——时，我们会赋予处理节点不同的地位。

处理节点与其他节点地位不同这一点，在第 2 章引入的因果树中也以图形方式明确体现：对应非处理变量 $L$ 与 $Y$ 的非处理分支被圈在圆圈中；第 5 章表示充分原因的“饼图”则区分了潜在处理 $A$ 与 $E$ 以及背景因素 $U$。此外，第 3 章关于处理的充分定义良好的干预的讨论，强调了加在处理变量 $A$ 上、而不适用于其他变量的要求。

相比之下，本章的因果图表面上赋予图中所有变量相同地位——当把因果图视为误差独立的非参数结构方程模型的表示时，情况确实如此（见技术注记 6.2）。因果图中所有变量表面上的平等地位可能造成误导，因为其中一些变量对应定义不良的干预。画一张因果图，把“肥胖”作为结局 $Y$ 甚至作为协变量 $L$ 的节点包含进去，也许是可以的（更多讨论见 9.5 节）。然而，出于第 3 章讨论的理由，把“肥胖”作为处理 $A$ 的节点包含进去，一般是不可以的。在因果图中，处理变量的节点必须对应充分定义良好的干预。

\mnote{近来 Pearl (2018, 2019) 提出了一种基于“听从”其他变量的变量的因果概念，它仍然假定每个变量都有定义良好的反事实。}

例如，假设我们关心“体重减轻” $A$ 对死亡率 $Y$ 的潜在因果效应，如第 3 章所讨论。图 6.10 的因果图包含 $A$ 与 $Y$ 的节点，也包含影响体重的因素的节点。为简化，因果图只包含其中 3 个因素：热量摄入 $Z$（我们假定它只能通过体重减轻影响死亡率）、运动 $L$（它可以通过体重减轻之外的途径影响死亡率），以及遗传特质 $U$（它可以通过同样独立于体重减轻的其他途径影响死亡率）。

\begin{figure}[htbp]
\centering
\begin{tikzpicture}[causal]
  \node[vtx] (f10L) at (0,2.0) {$L$};
  \node[vtx] (f10Z) at (0,0.65) {$Z$};
  \node[vtx] (f10U) at (0,-0.7) {$U$};
  \node[vtx] (f10A) at (2.5,0.65) {$A$};
  \node[vtx] (f10Y) at (5.2,0.65) {$Y$};
  \draw[ed] (f10L) -- (f10A);
  \draw[ed] (f10L) -- (f10Y);
  \draw[ed] (f10Z) -- (f10A);
  \draw[ed] (f10U) -- (f10A);
  \draw[ed] (f10U) -- (f10Y);
  \draw[ed] (f10A) -- (f10Y);
\end{tikzpicture}
\caption{}
\label{fig:6.10}
\end{figure}

识别并解释处理 $A$ 对结局 $Y$ 的效应，需要关于如何干预 $A$ 的知识。当存在若干种干预 $A$ 的可能方式，且其中一些可能干预对结局 $Y$ 有直接效应（如图 6.10）时，“$A$ 对 $Y$ 的效应”是什么意思就变得不清楚了。在我们的例子中，通过热量限制 $Z$ 减轻体重，与通过增加运动 $L$ 或通过基因操作 $U$ 减轻体重，会导致不同的死亡风险。即便有人愿意无视这个定义不良的因果效应，识别达到可交换性所需的变量也会是一项艰巨的挑战，如第 3 章所讨论。

明确说出所关心的干预，是获得定义良好的因果效应、识别相关数据、选择调整变量的重要一步。

\begin{finepoint}[因果结构的发现]
在本书中，我们把因果图用作表示我们关于手头问题因果结构的专家知识——或假设——的方式。也就是说，因果图指导数据分析。反过来如何？我们能否在不作因果结构假设的情况下通过数据分析来学习因果结构？通过数据分析学习因果结构各组成部分的过程称为\emph{发现}（discovery）。关于因果发现各种方法的描述，见 Spirtes 等 (2000) 与 Peters 等 (2017) 的书。

下面简要讨论在如下假设下的因果发现：观测数据来自一个未知的因果 DAG，该 DAG 除观测变量外还包含未知数目的未观测变量 $U$。若假定忠实性，因果发现有时是可能的，因为观测数据分布中的统计独立关系蕴含 DAG 上缺失的因果箭头。即便假定忠实性，发现也常常不可能。例如，假设我们在数据中发现两个变量 $B$ 与 $C$ 之间有强关联。我们无法得知涉及 $B$ 与 $C$ 的因果结构，因为它们的关联与许多因果图相容：$B$ 引起 $C$（$B\to C$）、$C$ 引起 $B$（$C\to B$）、$B$ 与 $C$ 共享未测量的原因 $U$（$B\leftarrow U\to C$）、$B$ 与 $C$ 有一个已被取条件的未观测共同效应 $U$，以及各种组合。若我们知道 $B$ 与 $C$ 的时间顺序，我们只能排除 $B\to C$（若 $C$ 早于 $B$）或 $C\to B$（若 $B$ 早于 $C$）的因果图。

然而，有些情境中从数据学习因果结构看来是可能的。例如，考虑 3 个变量 $Z$、$A$、$Y$，且我们知道其时间顺序为 $Z$ 最先、$A$ 其次、$Y$ 最后。在缺乏先验知识与经验数据时，图 6.11 表示一个可能的因果 DAG，因为我们无法排除 $Z$ 对 $A$ 与 $Y$ 有直接效应、$A$ 对 $Y$ 有效应，或任何一对变量之间存在未测量的共同原因。我们也无法排除删去一条或多条箭头的任何子图。

现在假设（假想地）我们获得了无穷多个个体在变量 $Z$、$A$、$Y$ 上的数据。数据分析发现 3 个变量两两边缘有关联，且唯一成立的条件独立是 $Z\indep Y\given A$。那么，若我们愿意假定忠实性成立，唯一与我们的分析相容的因果 DAG 是 $Z\to A\to Y$，或许再加上 $Z$ 与 $A$ 的共同原因 $U_1$（作为 $Z\to A$ 箭头的补充或替代）。这是因为，若 $Z$ 是 $Y$ 的父节点、或与 $Y$ 共享原因 $U_2$，或存在 $A$ 与 $Y$ 的未测量共同原因 $U_3$，则 $Z$ 与 $Y$ 不可能在给定 $A$ 时统计独立（假定忠实性）。因此，要解释 $Y$ 与 $A$ 的边缘相关，必须存在从 $A$ 到 $Y$ 的因果箭头。（注意，若存在 $A$ 与 $Y$ 的未测量共同原因，则不会发现任何条件独立，此时即便假定忠实性，我们也无法确定 $A$ 是否引起 $Y$。）

总之，所学习到的因果 DAG 蕴含 $Z$ 不是 $Y$ 的直接原因（父节点）、不存在 $A$ 与 $Y$ 的未测量共同原因，并且事实上 $A$ 对 $Y$ 的平均因果效应可由 $\E[Y\given A=1]-\E[Y\given A=0]$ 识别，因为可交换性成立。当然，我们并没有无穷大的样本量。关于有限样本对因果发现的含义的讨论，推迟到技术注记 10.7。

\begin{center}
\begin{tikzpicture}[causal]
  \node[vtx] (f11Z) at (0,1.1) {$Z$};
  \node[vtx] (f11A) at (2.7,1.1) {$A$};
  \node[vtx] (f11Y) at (5.4,1.1) {$Y$};
  \node[vtx] (f11U1) at (1.35,-0.25) {$U_1$};
  \node[vtx] (f11U3) at (4.05,-0.25) {$U_3$};
  \node[vtx] (f11U2) at (2.7,-1.7) {$U_2$};
  \draw[ed] (f11Z) -- (f11A);
  \draw[ed] (f11A) -- (f11Y);
  \draw[ed] (f11Z) to[bend left=32] (f11Y);
  \draw[ed] (f11U1) -- (f11Z);
  \draw[ed] (f11U1) -- (f11A);
  \draw[ed] (f11U3) -- (f11A);
  \draw[ed] (f11U3) -- (f11Y);
  \draw[ed] (f11U2) to[out=168,in=-78] (f11Z);
  \draw[ed] (f11U2) to[out=12,in=-102] (f11Y);
\end{tikzpicture}
\captionof{figure}{}
\label{fig:6.11}
\end{center}
\end{finepoint}

\section{偏倚的结构分类}

“偏倚”一词常被做因果推断的研究者使用。该词有若干相关但在技术上不同的用法（见第 10 章）。当数据不足以识别——计算——因果效应，即便样本量无穷大时，我们称存在\emph{系统性偏倚}。（本章由于假定样本量无穷大，偏倚指系统性偏倚。）非正式地说，我们常把系统性偏倚称为处理与结局之间任何并非来自目标总体中处理对结局之因果效应的结构性关联。由于因果图有助于表示不同的关联来源，我们可以用因果图按来源对系统性偏倚分类，从而深化关于偏倚的讨论。

\mnote{存在系统性偏倚时，没有任何估计量可以是一致的。一致估计量的定义见第 1 章。}

以我们此前各章讨论过的关键偏倚来源为例：受治疗者与未受治疗者之间缺乏可交换性。对于整个总体中的平均因果效应，当
\[
\Pr[Y^{a=1}=1]-\Pr[Y^{a=0}=1]\neq \Pr[Y=1\given A=1]-\Pr[Y=1\given A=0]
\]
时，我们说存在（无条件）偏倚；当（无条件）可交换性 $Y^{a}\indep A$ 不成立时即为这种情形。无（无条件）偏倚意味着总体中的关联度量（如关联风险比或风险差）是总体中相应效应度量（如因果风险比或风险差）的一致估计。

即便“处理对结局无因果效应”这一零假设成立，缺乏可交换性也会导致偏倚。也就是说，即便处理对结局没有因果效应，处理与结局在数据中仍会有关联。我们于是说，缺乏可交换性导致\emph{零假设下的偏倚}。在表 3.1 所概括的观察性研究中，存在零假设下的偏倚，因为因果风险比为 1，而关联风险比为 1.26。任何导致零假设下偏倚的因果结构，也会导致备择假设下（即处理对结局确有非零效应时）的偏倚。反之则不然。

\mnote{例如，以某些变量为条件可能在备择假设下（即偏离零假设时）造成选择偏倚，而在零假设下不会，如 Greenland (1977) 与 Hernán (2017) 所述。另见第 18 章。}

对于 $L$ 各水平内的平均因果效应，只要对至少一个层 $l$，
\[
\begin{aligned}
&\Pr[Y^{a=1}=1\given L=l]-\Pr[Y^{a=0}=1\given L=l]\\
&\qquad\neq \Pr[Y=1\given L=l,A=1]-\Pr[Y=1\given L=l,A=0],
\end{aligned}
\]
我们就说存在\emph{条件偏倚}；当条件可交换性 $Y^{a}\indep A\given L=l$ 并非对所有 $a$ 与 $l$ 成立时，通常就是这种情形。

本书迄今已多次提到缺乏可交换性。然而我们尚未探究产生缺乏可交换性的因果结构。有了因果图这一方法工具，我们就能描述缺乏可交换性如何由两种不同的因果结构产生：

\begin{enumerate}
\item \emph{共同原因}：当处理与结局共享共同原因时，关联度量一般不同于效应度量。许多流行病学家用“混杂”（confounding）一词指称这种偏倚。
\item \emph{以共同效应为条件}：这种结构是许多流行病学家称为零假设下选择偏倚的偏倚来源。
\end{enumerate}

第 7 章将聚焦于共同原因存在所导致的混杂偏倚，第 8 章聚焦于以共同效应为条件所导致的选择偏倚。再次说明，二者都是缺乏可交换性所致零假设下偏倚的例子。第 9 章将聚焦另一偏倚来源：测量误差。迄今我们假定所有变量——处理 $A$、结局 $Y$ 与协变量 $L$——都被完美测量。然而在实践中，某种程度的测量误差是意料之中的。测量误差导致的偏倚称为测量偏倚或信息偏倚。我们将看到，某些类型的测量偏倚也会导致零假设下的偏倚。

因此，随后三章我们把注意力转向三类系统性偏倚——混杂、选择与测量。这些偏倚在观察性研究与随机化实验中都可能出现。随机化实验对偏倚的易感性或许从前几章看不出来，因为在那里我们把观察性研究概念化为某种不完善的随机化实验，而只考虑理想的随机化实验：随访中无人失访、所有参与者都依从其被分配的处理、研究参与者与研究者都不知道处理分配。虽然我们对随机化实验的这种准神话式刻画有助于教学，真实的随机化实验却很少如此。第一部分余下各章将详述实验与观察之间那条有时模糊的界线。

\mnote{另一类偏倚也可能来自（非结构性的）随机变异。见第 10 章。}

在此之前，我们稍作迂回，描述存在效应修饰时的因果图。

\section{效应修饰的结构}

识别偏倚的潜在来源是因果图的一个关键用途：我们可以用因果专家知识画图，然后搜寻处理与结局之间关联的来源。因果图在说明第 4 章讨论的效应修饰概念方面帮助较小。

\mnote{用因果图对效应修饰作更细致的分类，见 VanderWeele 与 Robins (2007b)。}

假设在一项实验中随机分配心脏移植 $A$，以识别 $A$ 对死亡 $Y$ 的平均因果效应。为简化，假定没有偏倚，因此图 6.2 足以表示这项研究。计算 $A$ 对 $Y$ 之风险的影响并无困难。由于关联即因果，关联风险差 $\Pr[Y=1\given A=1]-\Pr[Y=1\given A=0]$ 可以解释为因果风险差 $\Pr[Y^{a=1}=1]-\Pr[Y^{a=0}=1]$。然而研究者想更进一步，因为他们怀疑心脏移植的因果效应随参与研究的各医院所提供的医疗质量而异。于是研究者把所有个体分为接受高质量（$V=1$）或常规质量（$V=0$）医疗者，按第 4 章所述计算 $V$ 各水平内的分层风险差，并确实证实存在 $V$ 在加法尺度上的效应修饰。图 6.12 的因果图包含效应修饰因子 $V$，有一条箭头指向结局 $Y$，但没有箭头指向处理 $A$（$A$ 被随机分配，因而独立于 $V$）。有两点需要提醒。

\begin{figure}[htbp]
\centering
\begin{tikzpicture}[causal]
  \node[vtx] (f12V) at (0,1.05) {$V$};
  \node[vtx] (f12A) at (2.5,0.15) {$A$};
  \node[vtx] (f12Y) at (5.1,0.15) {$Y$};
  \draw[ed] (f12V) to[bend left=26] (f12Y);
  \draw[ed] (f12A) -- (f12Y);
\end{tikzpicture}
\caption{}
\label{fig:6.12}
\end{figure}

第一，图 6.12 的因果图即便不包含 $V$ 也仍是有效的因果图，因为 $V$ 不是 $A$ 与 $Y$ 的共同原因。只是因为因果问题提到了 $V$（即 $A$ 对 $Y$ 的平均因果效应在 $V$ 各水平内是多少？），才需要把 $V$ 纳入因果图。在“医疗质量” $V$ 与结局 $Y$ 之间路径上测得的其他变量也可能有资格作为效应修饰因子。例如，图 6.13 显示效应修饰因子“治疗并发症” $N$，它部分中介了 $V$ 对 $Y$ 的效应。

\begin{figure}[htbp]
\centering
\begin{tikzpicture}[causal]
  \node[vtx] (f13V) at (0,1.0) {$V$};
  \node[vtx] (f13N) at (1.7,2.0) {$N$};
  \node[vtx] (f13A) at (2.5,0.15) {$A$};
  \node[vtx] (f13Y) at (5.1,0.15) {$Y$};
  \draw[ed] (f13V) -- (f13N);
  \draw[ed] (f13N) -- (f13Y);
  \draw[ed] (f13V) to[bend left=24] (f13Y);
  \draw[ed] (f13A) -- (f13Y);
\end{tikzpicture}
\caption{}
\label{fig:6.13}
\end{figure}

第二，图 6.12 的因果图并不必然表示存在 $V$ 的效应修饰。该因果图蕴含 $A$ 与 $V$ 都影响死亡 $Y$，但它并不区分 $V$ 可能修饰 $A$ 对 $Y$ 之效应的以下三种定性上不同的方式：

\begin{enumerate}
\item 处理 $A$ 对死亡率 $Y$ 的因果效应在 $V=1$ 层与 $V=0$ 层方向相同（有害或有益）。
\item 处理 $A$ 对死亡率 $Y$ 的因果效应在 $V=1$ 层的方向与 $V=0$ 层相反，即存在\emph{定性效应修饰}。
\item 处理 $A$ 在 $V$ 的某一层对 $Y$ 有因果效应，而在另一层没有因果效应，例如 $A$ 只杀死 $V=0$ 的个体。
\end{enumerate}

也就是说，像图 6.12 这样的有效因果图无法区分上述三种定性上不同的 $V$ 的效应修饰类型。

在上例中，效应修饰因子 $V$ 对结局有因果效应。然而许多效应修饰因子对结局并没有因果效应，它们只是对结局有因果效应的变量的替代物。图 6.14 包含变量“治疗费用” $S$（1：高，0：低），它受“医疗质量” $V$ 影响，但本身对死亡率 $Y$ 没有效应。按 $S$（而非按 $V$）分层的分析一般会检出 $S$ 的效应修饰，尽管真正修饰 $A$ 对 $Y$ 之效应的是变量 $V$。变量 $S$ 是\emph{替代效应修饰因子}，而变量 $V$ 是\emph{因果效应修饰因子}（见 4.2 节）。由于因果效应修饰因子与替代效应修饰因子在实践中常常难以区分，效应修饰这一概念兼指二者。如 4.2 节所讨论，有些人偏好使用中性的说法“因果效应的异质性”而非“效应修饰”，以避免混淆。例如，有人可能忍不住把“费用修饰了心脏移植对死亡率的效应，因为费用越高效应越有益”这句话解释为应当提高医疗价格而不必提高其质量的理由。

\begin{figure}[htbp]
\centering
\begin{tikzpicture}[causal]
  \node[vtx] (f14V) at (0,1.0) {$V$};
  \node[vtx] (f14S) at (1.7,2.0) {$S$};
  \node[vtx] (f14A) at (2.5,0.15) {$A$};
  \node[vtx] (f14Y) at (5.1,0.15) {$Y$};
  \draw[ed] (f14V) -- (f14S);
  \draw[ed] (f14V) to[bend left=24] (f14Y);
  \draw[ed] (f14A) -- (f14Y);
\end{tikzpicture}
\caption{}
\label{fig:6.14}
\end{figure}

替代效应修饰因子只是与因果效应修饰因子有关联的变量。图 6.14 描绘的是这种关联来自因果效应修饰因子对替代效应修饰因子之效应的情境。然而这种关联也可能来自共享共同原因或以共同效应为条件。例如，图 6.15 包含变量“居住地”（1：希腊，0：罗马）$U$ 与“护照所定国籍” $P$（1：希腊，0：罗马）。居住地 $U$ 是医疗质量 $V$ 与国籍 $P$ 的共同原因。因此 $P$ 会表现为替代效应修饰因子，因为 $P$ 与因果效应修饰因子 $V$ 有关联。另一个（确实有点荒唐的）说明这一问题的例子：图 6.16 包含变量“医疗费用” $S$ 与“在医院饮用瓶装矿泉水（而非自来水）” $W$。在发达国家，使用矿泉水 $W$ 影响费用 $S$，但不影响死亡率 $Y$。若研究限制在低费用医院（$S=0$），那么使用矿泉水 $W$ 一般会与医疗质量 $V$ 有关联，于是 $W$ 会表现为替代效应修饰因子。总之，替代效应修饰因子可以通过包括共同原因、以共同效应为条件或原因与结果在内的各种结构与因果效应修饰因子相关联。

\mnote{关于低费用医院中 $W$ 与 $V$ 关联的直观解释：使用矿泉水的低费用医院必须通过少花钱于能降低死亡率的医疗环节来抵消矿泉水的费用。于是，在低费用医院中，矿泉水的使用与医疗质量呈反向关联。}

\begin{figure}[htbp]
\centering
\begin{tikzpicture}[causal]
  \node[vtx] (f15U) at (0,1.75) {$U$};
  \node[vtx] (f15P) at (1.7,2.35) {$P$};
  \node[vtx] (f15V) at (0.3,0.15) {$V$};
  \node[vtx] (f15A) at (2.6,0.15) {$A$};
  \node[vtx] (f15Y) at (5.1,0.15) {$Y$};
  \draw[ed] (f15U) -- (f15P);
  \draw[ed] (f15U) -- (f15V);
  \draw[ed] (f15V) to[bend left=24] (f15Y);
  \draw[ed] (f15A) -- (f15Y);
\end{tikzpicture}
\caption{}
\label{fig:6.15}
\end{figure}

\begin{figure}[htbp]
\centering
\begin{tikzpicture}[causal]
  \node[vtx] (f16W) at (0,2.45) {$W$};
  \node[vtx] (f16S) at (1.7,1.55) {$S$};
  \node[vtx] (f16V) at (0,0.15) {$V$};
  \node[vtx] (f16A) at (2.6,0.15) {$A$};
  \node[vtx] (f16Y) at (5.1,0.15) {$Y$};
  \draw[ed] (f16W) -- (f16S);
  \draw[ed] (f16V) -- (f16S);
  \draw[ed] (f16V) to[bend left=24] (f16Y);
  \draw[ed] (f16A) -- (f16Y);
  \node[draw,inner sep=2.6pt,line width=0.5pt,fit=(f16S)] {};
\end{tikzpicture}
\caption{}
\label{fig:6.16}
\end{figure}

因果图原则上对两种处理 $A$ 与 $E$ 之间是否存在交互作用不作断言。然而，当因果图用表示充分成分原因的节点（见第 5 章）加以扩充，即加入从处理指向充分成分原因的确定性箭头时，它就能编码关于交互作用的信息。由于交互作用的存在会影响以共同效应为条件所致关联的大小与方向，这些扩充后的因果图将在第 8 章讨论。

\begin{finepoint}[再论非干预变量的因果效应]
如 3.4 节所承诺，我们用因果 DAG 讨论可能使一些研究者认为非干预变量（如血压）具有定义良好的反事实的那类经验证据。

考虑一项研究，考察药物 $Z$ 对收缩压 $A$（一种状态）及随后中风风险 $Y$（二值结局）的效应。假设 (i) $Z$ 与 $A$ 和 $Y$ 都有关联，(ii) $A$ 与 $Y$ 有关联，(iii) 给定 $A$ 时 $Z$ 与 $Y$ 独立。那么，如精要 6.2 所示，唯一忠实的因果 DAG 有 $A\to Y$ 箭头，没有从 $Z$ 到 $Y$ 的箭头，也没有 $A$ 与 $Y$ 或 $Z$ 与 $Y$ 的未测量共同原因。若反事实 $Y^{a}$ 定义良好，我们可以在忠实性下断定：(a) 不存在 $A$ 与 $Y$ 的其他未测量共同原因 $U$，(b) 可交换性 $Y^{a}\indep A$ 成立，(c) $A$ 对 $Y$ 的平均因果效应等于 $\E[Y\given A=1]-\E[Y\given A=0]$，且 (d) 图 6.11 中的 $Z$ 或 $U_1$（或二者）只通过 $A$ 影响 $Y$。

假设在研究之前，某研究者因 $A$ 无法被直接干预而不认为反事实 $Y^{a}$——从而是 $A$ 对 $Y$ 的因果效应——定义良好。这位研究者现在可能把 (i)--(iii) 视为强经验证据，表明 $Y^{a}$ 定义良好且 (a)--(d) 为真，因为在任何忠实的因果 DAG 上，$Z$ 或 $U_1$ 中的一个或两个只通过 $A$ 影响 $Y$。若 (i)--(iii) 只在给定基线协变量 $L$ 时成立，则 (a) 与 (d) 不变，(b) 换为条件可交换性 $Y^{a}\indep A\given L$ 成立，(c) 换为在任何 $L=l$ 内 $A$ 对 $Y$ 的平均因果效应等于 $\E[Y\given A=1,L=l]-\E[Y\given A=0,L=l]$。

遗憾的是，在多数情形中仍会存在 $A$ 与 $Y$ 的未测量（且常常未知的）共同原因 $U_3$，因此 (iii) 不成立，因为给定对撞节点 $A$ 时 $Z$ 与 $Y$ 会有关联。在这种情况下，即便假定忠实性，图 6.11 中的 $Z\to Y$ 箭头也不能再删去。结果是，数据不会为“$Y^{a}$ 定义良好”提供任何经验支持，因为我们不再能排除 $Z$ 对 $Y$ 有不经 $A$ 的直接效应。

这就提出一个问题：当存在 $A$ 与 $Y$ 的未测量共同原因 $U_3$ 时，人们究竟能否获得“状态 $A$ 具有定义良好的反事实结局 $Y^{a}$”的经验证据？下面我们说明，在忠实性下答案是肯定的：办法是做一个四臂析因剂量探索随机化试验，其研究设计类似于常用来研究抗高血压药在动物与人体中效应的设计。

在我们的剂量探索试验中，血压值相近的个体被随机分配到两种降压药之一（$Z=1$ 或 $Z=2$），并随机分配到两种血压降低值之一（$A=a_1$ 或 $A=a_2$，单位 mmHg）。在试验最初几天，对每个个体分别逐步增加所分配药物 $Z$ 的日剂量 $D$，直到该个体的平均血压达到其被分配的取值 $A$。我们假定这是一个理想的剂量探索试验：每个个体都达到了其被分配的降压值 $A$，剂量 $D$ 在整个试验期间持续维持每个个体被分配的 $A$，且短暂的剂量探索期对结局 $Y$ 没有影响。

假设在我们的理想剂量探索试验中观测到结局 $Y$ 在给定 $A$ 时独立于 $D$，即 $\E[Y\given D=2,A]=\E[Y\given D=1,A]=\E[Y\given A]$，且 $\E[Y\given A=a_1]\neq\E[Y\given A=a_2]$。那么，在忠实性下，并假定 $Y^{d,a}$ 与 $Y^{a}$ 定义良好，我们可以断定 $D$ 对 $Y$ 没有不经 $A$ 的直接效应，且平均因果效应 $\E[Y^{a_1}-Y^{a_2}]$ 可由 $\E[Y\given A=a_1]-\E[Y\given A=a_2]$ 识别，无论是否存在 $A$ 与 $Y$ 的未知共同原因 $U$。为解释原因，我们比较剂量探索随机化试验与标准（非剂量探索）试验——后者随机分配处理 $Z$ 与剂量 $D$——所对应的因果 DAG。

图 6.17 显示标准试验的因果 DAG，其中 $U$ 是 $A$ 与 $Y$ 的未测量共同原因。由于给定 $A$ 时 $Y$ 与 $Z$ 有关联，我们无法凭经验判定虚线箭头是否存在，因此无法区分 $Z$ 对 $Y$ 不经 $A$ 的直接效应与 $Z$ 对 $Y$ 只通过 $A$ 的直接效应。图 6.18 显示我们剂量探索试验的因果 DAG，其中观测到的血压 $A$ 等于（从而是其确定性函数）被分配的 $A_{\mathrm{ass}}$。因此在剂量探索试验中，除来自 $A_{\mathrm{ass}}$ 的箭头外，不可能有指向 $A$ 的箭头。达到 $A=A_{\mathrm{ass}}$ 所需的剂量 $D$ 由 $A_{\mathrm{ass}}$、$Z$ 与 $U$ 决定。在忠实性下，给定 $A$ 时 $Z$ 独立于 $Y$，当且仅当 $Z$ 对 $Y$ 没有不经 $A$ 的直接效应。此外，由于 $\E[Y\given A=a_1]\neq\E[Y\given A=a_2]$，必存在 $A\to Y$ 箭头，从而 $A$ 对 $Y$ 的效应由 $\E[Y\given A=a_1]-\E[Y\given A=a_2]$ 给出。

\begin{center}
\begin{tikzpicture}[causal]
  \node[vtx] (f17Z) at (0,1.7) {$Z$};
  \node[vtx] (f17A) at (2.6,1.7) {$A$};
  \node[vtx] (f17Y) at (5.3,1.7) {$Y$};
  \node[vtx] (f17D) at (0.5,0.25) {$D$};
  \node[vtx] (f17U) at (3.6,0.25) {$U$};
  \draw[ed] (f17Z) -- (f17A);
  \draw[ed] (f17A) -- (f17Y);
  \draw[ed] (f17D) -- (f17A);
  \draw[ed] (f17U) -- (f17A);
  \draw[ed] (f17U) -- (f17Y);
  \draw[edd] (f17Z) to[bend left=24] (f17Y);
  \draw[edd] (f17D) to[bend right=24] (f17Y);
\end{tikzpicture}
\captionof{figure}{}
\label{fig:6.17}
\end{center}

\begin{center}
\begin{tikzpicture}[causal]
  \node[vtx] (f18R) at (0,2.5) {$A_{\mathrm{ass}}$};
  \node[vtx] (f18A) at (3.0,2.5) {$A$};
  \node[vtx] (f18Y) at (5.9,1.55) {$Y$};
  \node[vtx] (f18D) at (2.3,0.85) {$D$};
  \node[vtx] (f18Z) at (0,0) {$Z$};
  \node[vtx] (f18U) at (4.2,0) {$U$};
  \draw[edb] (f18R) -- (f18A);
  \draw[ed] (f18R) -- (f18D);
  \draw[ed] (f18Z) -- (f18D);
  \draw[ed] (f18U) -- (f18D);
  \draw[edd] (f18D) -- (f18Y);
  \draw[ed] (f18A) -- (f18Y);
  \draw[ed] (f18U) -- (f18Y);
  \draw[edd] (f18Z) to[bend right=26] (f18Y);
\end{tikzpicture}
\captionof{figure}{}
\label{fig:6.18}
\end{center}

如上所述，假设某研究者因 $A$ 无法被直接干预而不认为反事实 $Y^{a}$——从而是 $A$ 对 $Y$ 的因果效应——定义良好。然而在忠实性下，这位研究者可以把 $A\to Y$ 箭头的必要性视为 $Y^{a}$ 定义良好的强经验证据。
\end{finepoint}
